% ---
% id: protocol-compiler/rem:sidecond
% kind: rem
% title: "Necessity of the key-rule side condition"
% summary: >-
%   Without the condition that the key rule map into the key values, the
%   specification's sifting predicate stops witnessing announcement-determined
%   sifting. Check V5 catches only a key rule that withholds a key on part of a
%   passing class, so the condition is a clause of [[A-SPEC]].
% prominence: standard
% status: accepted
% uses: [def:spec]
% mentions: [thm:sifting-and-key, def:declaration, def:validator, ass:spec]
% assumes: []
% refs: []
% provenance: newly-derived
% lean:
% ---

\begin{remark}\label{rem:sidecond}
Drop the requirement that $r^{0}_{\hold}$ map into $\Zval$, and
Theorem~\ref{thm:sifting-and-key}\,(ii) fails: a key rule that returns $\perp$
on a \emph{passing} class makes $r(x,v)=\perp$ for pairs whose announcement
class $\varsigma$ calls $\mathrm{pass}$, so the specification's $\varsigma$ no
longer witnesses announcement-determined sifting. The two ways this can happen
differ in what is still decided. Where the key rule returns $\perp$ on part of
a passing class, no predicate witnesses D(v): whether a round is sifted is not
determined by the announcements, a sifting event becomes invisible to the
public record, and check V5 of Definition~\ref{def:validator} refuses the
candidate. Where it returns $\perp$ on the whole of a passing class, D(v) still
holds, witnessed by a predicate that calls that class $\mathrm{fail}$, and V5
passes; what is lost is that the specification's $\varsigma$ is that witness,
and \eqref{eq:chkey} of Definition~\ref{def:declaration}, which reads the
specification's $\varsigma$, may then count as key-live an acceptance value
that no key-producing class reaches. This is why the condition is stated in
the specification's type (Definition~\ref{def:spec}\,(d)) and carried as a
clause of Assumption~\ref{ass:spec}, V5 catching only its first failure.
\end{remark}
